\documentclass[a4paper,10pt]{article}
\usepackage[utf8]{inputenc}
\usepackage[T1]{fontenc}
\usepackage[spanish,es-tabla]{babel}
\usepackage{geometry}
\geometry{top=0.8cm, bottom=0.8cm, left=1.2cm, right=1.2cm}
\usepackage{xcolor}
\usepackage{hyperref}
\usepackage{enumitem}
\usepackage{titlesec}
\usepackage{tikz}

% Soporte para caracteres Unicode subíndices (SpO₂)
\DeclareUnicodeCharacter{2082}{\textsubscript{2}}

% Paleta ejecutiva Navy & Slate
\definecolor{primary}{RGB}{15, 34, 64}        % Deep Midnight Navy
\definecolor{rulecolor}{RGB}{203, 213, 225}    % Muted Slate Rule (#cbd5e1)
\definecolor{textdark}{RGB}{30, 41, 59}        % Slate 800 (#1e293b)
\definecolor{subtext}{RGB}{100, 116, 139}      % Slate 500 (#64748b)
\definecolor{accentlink}{RGB}{20, 50, 95}     % Deep Navy link

\IfFileExists{inter.sty}{
  \usepackage[default]{inter}
}{
  \usepackage{tgheros}
  \renewcommand{\familydefault}{\sfdefault}
}

\IfFileExists{fontawesome5.sty}{
  \usepackage{fontawesome5}
}{
  \newcommand{\faEnvelope}{\tikz[baseline=-0.3ex]{\draw[line width=0.55pt, rounded corners=0.5pt, draw=subtext] (0,0) rectangle (0.32,0.22); \draw[line width=0.45pt, draw=subtext] (0,0.22) -- (0.16,0.10) -- (0.32,0.22);}}
  \newcommand{\faGlobe}{\tikz[baseline=-0.3ex]{\draw[line width=0.55pt, draw=subtext] (0,0) circle (0.11); \draw[line width=0.35pt, draw=subtext] (-0.11,0) -- (0.11,0); \draw[line width=0.35pt, draw=subtext] (0,-0.11) -- (0,0.11); \draw[line width=0.35pt, draw=subtext] (0,0) ellipse (0.05 and 0.11);}}
  \newcommand{\faLinkedin}{\tikz[baseline=-0.4ex]{\node[fill=primary, text=white, rounded corners=0.8pt, inner sep=0.8pt, font=\fontsize{3.5}{4.5}\selectfont\bfseries] {in};}}
  \expandafter\newcommand\csname faMapMarker*\endcsname{\tikz[baseline=-0.4ex]{\draw[line width=0.5pt, fill=subtext!50, draw=subtext] (0,0.06) circle (0.05); \draw[line width=0.45pt, draw=subtext] (-0.04,0.03) -- (0,-0.07) -- (0.04,0.03);}}
  \expandafter\newcommand\csname faExternalLink*\endcsname{\tikz[baseline=-0.2ex]{\draw[line width=0.45pt, draw=accentlink] (0,0) -- (0.14,0) -- (0.14,0.06); \draw[line width=0.45pt, draw=accentlink] (0,0) -- (0.14,0) -- (0.06,0.14); \draw[line width=0.45pt, draw=accentlink] (0.04,0.04) -- (0.17,0.17); \draw[line width=0.45pt, draw=accentlink] (0.10,0.17) -- (0.17,0.17) -- (0.17,0.10);}}
}

\hypersetup{
  colorlinks=true,
  urlcolor=accentlink,
  linkcolor=accentlink
}

\titleformat{\section}{\large\bfseries\color{primary}}{}{0em}{}[\vspace{1.5pt}\color{rulecolor}\titlerule]
\titlespacing*{\section}{0pt}{4.5pt}{2pt}

\setlist[itemize]{label=\textcolor{primary}{\small\textbullet}, leftmargin=1.15em, itemsep=1.0pt, topsep=1.0pt, parsep=0pt}

\linespread{1.0}
\pagestyle{empty}

\begin{document}
\color{textdark}

% Cabecera centrada
\begin{center}
  {\Huge\bfseries\color{primary} Diego De Pablo} \par\vspace{2.5pt}
  {\large\textbf{\color{textdark}Ingeniero de Software \textbar\ Backend, Datos e IA Aplicada}} \par\vspace{3.5pt}
  {\small\color{subtext}
    \href{mailto:diegodepablo.programa@gmail.com}{\faEnvelope\ diegodepablo.programa@gmail.com} \enskip $\cdot$ \enskip
    \href{https://linkedin.com/in/diego-de-pablo/}{\faLinkedin\ linkedin.com/in/diego-de-pablo} \enskip $\cdot$ \enskip
    \href{https://diegodepablo.is-a.dev/}{\faGlobe\ diegodepablo.is-a.dev} \enskip $\cdot$ \enskip
    {\csname faMapMarker*\endcsname\ Málaga, España}
  }
\end{center}

\vspace{-3pt}

\section*{Perfil Profesional}
Ingeniero de software y graduado en Bioinformática centrado en desarrollo backend, ingeniería de datos e IA aplicada. Experiencia construyendo servicios web modulares, pipelines de datos concurrentes y motores de búsqueda semántica con integración de modelos. Graduado con honores académicos y actualmente cursando el Máster Universitario en Ingeniería Informática.

\section*{Experiencia Profesional}

\noindent\textbf{Software Engineer, Khaos Research} \hfill {\color{subtext}\textbf{Sept 2025 -- Actualidad}}
\begin{itemize}
  \item Desarrollé servicios backend y módulos de integración para el proyecto europeo \textbf{SEDIA}, implementando especificaciones de espacios de datos federados e interoperables (Gaia-X).
  \item Diseñé e integré un motor de búsqueda híbrida combinando recuperación vectorial en \textbf{Qdrant} y búsqueda léxica BM25 en \textbf{Elasticsearch}, estructurado mediante agentes validados con \textbf{Pydantic AI}.
  \item Implementé el pipeline concurrente \textbf{MetaDataXtract} para la ingesta y extracción de metadatos en +30 formatos científicos, verificación de integridad SHA-256 y consulta automatizada a la API de VirusTotal para análisis de seguridad.
\end{itemize}

\vspace{2.5pt}
\noindent\textbf{Desarrollador Full-Stack, mSurgery} \hfill {\color{subtext}\textbf{Marzo 2025 -- Julio 2025}}
\begin{itemize}
  \item Implementé la migración del frontend hacia \textbf{Svelte 5} adoptando runes, desacoplando la gestión del estado para optimizar la reactividad del visor quirúrgico 3D.
  \item Desarrollé interfaces interactivas para la plataforma médica \textbf{SMOTTS}, integrando telemetría de gemelos digitales en tiempo real para simulación de procedimientos.
  \item Realicé demostraciones técnicas de la plataforma en inglés ante inversores internacionales y comités evaluadores en el Polo Nacional de Contenidos Digitales.
\end{itemize}

\vspace{2.5pt}
\noindent\textbf{Ingeniero de Software Júnior, T\&B Company Group} \hfill {\color{subtext}\textbf{Junio 2021 -- Diciembre 2021}}
\begin{itemize}
  \item Refactoricé módulos heredados en \textbf{Java} y \textbf{Angular} hacia arquitecturas desacopladas, resolviendo fallos en producción y mejorando la mantenibilidad del código.
  \item Diseñé contratos de interfaz estandarizados con \textbf{OpenAPI/Swagger} y documenté microservicios REST para agilizar la integración entre equipos.
\end{itemize}

\section*{Proyectos Técnicos Seleccionados}

\noindent\textbf{\href{https://diegodepablo.is-a.dev/projects/metadatasearch}{\csname faExternalLink*\endcsname\ MetaDataSearch: Buscador Semántico con IA}}
\\[1pt]
{\footnotesize\color{subtext}\textbf{Stack:} Next.js, FastAPI, Qdrant, Elasticsearch, Pydantic AI, Docker}
\begin{itemize}
  \item Sistema de búsqueda híbrida para catálogos científicos: combina vectorización densa y léxica, generando filtros semánticos estructurados a partir de consultas en lenguaje natural.
\end{itemize}

\vspace{2.5pt}
\noindent\textbf{\href{https://diegodepablo.is-a.dev/projects/tfg-patient-monitoring}{\csname faExternalLink*\endcsname\ Pulsera IoT de Telemonitorización Biomédica (TFG)}}
\\[1pt]
{\footnotesize\color{subtext}\textbf{Stack:} C++, FreeRTOS, ESP32, MAX30102, MPU-6050, MQTT, FastAPI, TimescaleDB}
\begin{itemize}
  \item Prototipo end-to-end de telemonitorización: firmware en ESP32 para adquisición de constantes vitales por PPG (pulso y SpO2) y detección de caídas con acelerometría, sincronización vía MQTT y panel clínico web. \textbf{Calificación 10/10 y Matrícula de Honor}.
\end{itemize}

\vspace{2.5pt}
\noindent\textbf{\href{https://diegodepablo.is-a.dev/projects/titan-workflow}{\csname faExternalLink*\endcsname\ TITAN Workflow: Asistente de Flujos Científicos}}
\\[1pt]
{\footnotesize\color{subtext}\textbf{Stack:} FastAPI, Pydantic AI, LangChain, Docker, LLM}
\begin{itemize}
  \item Servicio de composición de flujos científicos asistido por IA: valida dependencias de datos y esquemas de entrada/salida mediante contratos estrictos en Pydantic AI previo a su planificación.
\end{itemize}

\section*{Formación Académica}

\noindent\textbf{Grado en Ingeniería de la Salud (Mención Bioinformática)} \textbar\ \textit{Universidad de Málaga} \hfill {\color{subtext}\textbf{2021 -- 2025}}
\begin{itemize}
  \item Matrículas de Honor en asignaturas clave: \textbf{Fundamentos de Programación}, \textbf{Sistemas Inteligentes} y \textbf{Proyectos de Bioinformática}.
  \item Trabajo Fin de Grado (TFG) en IoT y Telemedicina calificado con \textbf{10/10 y Matrícula de Honor}.
\end{itemize}

\vspace{2.5pt}
\noindent\textbf{Máster Universitario en Ingeniería Informática} \textbar\ \textit{Universidad de Málaga} \hfill {\color{subtext}\textbf{2025 -- Actualidad}}
\begin{itemize}
  \item Formación avanzada en sistemas distribuidos, arquitecturas de datos e inteligencia artificial aplicada.
\end{itemize}

\section*{Competencias Técnicas e Idiomas}

\noindent\textbf{Lenguajes:} Python, TypeScript, JavaScript, C++, SQL, Java, Bash, R.\\[2pt]
\textbf{Backend y Datos:} FastAPI, PostgreSQL, Elasticsearch, TimescaleDB, Qdrant (Vector DB), REST APIs, OpenAPI/Swagger.\\[2pt]
\textbf{IA y Aprendizaje Automático:} Pydantic AI, LangChain, LLMs / Embeddings, PyTorch (académico y proyectos).\\[2pt]
\textbf{DevOps, Sistemas y Web:} Docker, Kubernetes, Argo CD, CI/CD (GitHub Actions), Linux, FreeRTOS (ESP32), React, Svelte 5.\\[2pt]
\textbf{Idiomas:} \textbf{Español} (Nativo) \enskip $\cdot$ \enskip \textbf{Inglés} (B2 Profesional / demostraciones técnicas internacionales).

\end{document}
